%% Generated by verification/make_figures.py from regular-graph-harmonic-quotient-examples.tex.
%% Figure fig:harmonic-quotient_d3reg12n+twinC5 (replaces figures/harmonic-quotient_d3reg12n+twinC5-G1.tex).
%% Gamma_1 = image of phi: 6 vertices, 7 edges.
%% Vertex y is filled with color hq<y>.
%% Layout: force-directed (0 crossings).
\begin{tikzpicture}[hqgraph]
\node[hqv, fill=hq0] (v0) at (-1.770,0.000) {0};
\node[hqv, fill=hq1] (v1) at (-0.457,-0.000) {1};
\node[hqv, fill=hq2] (v2) at (-0.259,-1.592) {2};
\node[hqv, fill=hq3] (v3) at (1.770,-1.030) {3};
\node[hqv, fill=hq4] (v4) at (1.770,1.030) {4};
\node[hqv, fill=hq5] (v5) at (-0.259,1.592) {5};
\draw[hqe] (v0) -- (v2);
\draw[hqe] (v0) -- (v5);
\draw[hqe] (v1) -- (v2);
\draw[hqe] (v1) -- (v5);
\draw[hqe] (v2) -- (v3);
\draw[hqe] (v3) -- (v4);
\draw[hqe] (v4) -- (v5);
\end{tikzpicture}%
